\newcommand{\hymnusmatutinum}{\pars{Hymnus.}

{
\grechangedim{interwordspacetext}{0.10 cm plus 0.15 cm minus 0.05 cm}{scalable}%
\antiphona{II}{temporalia/hym-GalliCantu.gtex}
\grechangedim{interwordspacetext}{0.22 cm plus 0.15 cm minus 0.05 cm}{scalable}%
}}
\newcommand{\lectioi}{\pars{Lectio I.} \scriptura{Prov. 15, 8-9. 16-17. 25-26. 29. 32}

\noindent De libro Proverbiórum.

\noindent Víctimæ impiórum abominábiles Dómino: vota iustórum grata sunt ei.

\noindent Abominátio est Dómino via ímpii; qui séquitur iustítiam, diligétur.

\noindent Mélius est parum cum timóre Dómini quam thesáuri magni cum sollicitúdine.

\noindent Mélius est deménsum ólerum cum caritáte quam vítulus saginátus cum ódio.

\noindent Domum superbórum demoliétur Dóminus et firmos fáciet términos víduæ.

\noindent Abominátio Dómini cogitatiónes malæ, et purus sermo pulchérrimus.

\noindent Longe est Dóminus ab ímpiis et oratiónes iustórum exáudiet.}
\newcommand{\responsoriumi}{\pars{Responsorium 1.} \scriptura{\Rbardot{} Prov. 1, 32-33 \Vbardot{} ibid. 8, 4; \textbf{H401}}

\vspace{-5mm}

\responsorium{VI}{temporalia/resp-aversioparvulorum-CROCHU.gtex}{}}
\newcommand{\lectioii}{\pars{Lectio II.} \scriptura{Prov. 15, 33; 16, 1-9; 17, 5}

\noindent Timor Dómini disciplína sapiéntiæ, et glóriam præcédit humílitas.

\noindent Hóminis est ánimum præparáre, et Dómini est respónsio linguæ.

\noindent Omnes viæ hóminis puræ sunt óculis eius, spirítuum ponderátor est Dóminus.

\noindent Revéla Dómino ópera tua, et dirigéntur cogitatiónes tuæ.

\noindent Univérsa secúndum próprium finem operátus est Dóminus; ímpium quoque ad diem malum.

\noindent Abominátio Dómini est omnis árrogans; etiámsi manus in manu fúerit, non erit ínnocens.

\noindent Misericórdia et veritáte redímitur iníquitas, et in timóre Dómini declinátur a malo.

\noindent Cum placúerint Dómino viæ hóminis, inimícos quoque eius convértet ad pacem.

\noindent Mélius est parum cum iustítia quam multi fructus sine æquitáte.

\noindent Cor hóminis dispónit viam suam, sed Dómini est dirígere gressus eius.

\noindent Qui déspicit páuperem, éxprobrat Factóri eius; et, qui in ruína lætátur alteríus, non erit impunítus.}
\newcommand{\responsoriumii}{\pars{Responsorium 2.} \scriptura{\Rbardot{} Prov. 1, 8 \Vbardot{} ibid. 3, 9; \textbf{H401}}

\vspace{-7mm}

\responsorium{VII}{temporalia/resp-audifilimihidisciplinam-CROCHU.gtex}{}}
\newcommand{\lectioiii}{\pars{Lectio III.} \scriptura{Nn. 51.53 : Funk 1,165-169)}

\noindent Ex Epístula sancti Cleméntis papæ Primi ad Corínthios.

\noindent Quæcúmque delíquimus et fécimus sedúcti a quodam ex adversárii servis, eórum remissiónem implorémus; illi vero qui seditiónis et dissensiónis príncipes exstitérunt, commúnem spem contemplári debent.

\noindent Qui enim in timóre et caritáte vitam degunt, se ipsos pótius quam próximos in cruciátus incídere volunt, et se ipsos pótius quam pulchre et iuste tráditam nobis consensiónem subíre vituperatiónem volunt.

\noindent Mélius enim est hómini peccáta sua confitéri, quam induráre cor suum; prout indurátum est eórum cor, qui contra Móysen fámulum Dómini seditiónem concitárunt, quorum condemnátio manifésta fuit.

\noindent \emph{Descendérunt enim in inférnum vivi,} et mors eos pascet. Phárao una cum exércitu suo et omnes præfécti Ægýpti, currus étiam et consisténtes in iis non áliam ob causam rubri maris profúndo immérsi sunt et periérunt, quam quod stulta eórum corda post tot signa et mirácula osténsa per fámulum Dei Móysen in terra Ægýpti induráta fúerint.

\noindent Quis ígitur inter vos generósus, quis miséricors, quis caritáte plenus?

\noindent Dicat: «Si propter me sedítio et discórdia et schísmata orta sunt, discédo, ábeo, quocúmque voluéritis, et quæ multitúdo iússerit, fácio; dúmmodo Christi grex in pace degat cum constitútis presbýteris.» Qui hoc fécerit, magnam sibi glóriam in Christo conciliábit, et omnis locus eum excípiet. Dómini enim est terra et plenitúdo eius. Hæc fáciunt et fácient, qui divínam, cuius numquam pǽnitet, vitam vivunt.}
\newcommand{\responsoriumiii}{\pars{Responsorium 3.} \scriptura{\Rbardot{} Prv. 23, 26 \Vbardot{} ibid. 5, 1; \textbf{H400}}

\vspace{-5mm}

\responsorium{I}{temporalia/resp-praebefili-CROCHU-cumdox.gtex}{}

\rubrica{vel ad libitum:}

\vspace{3mm}

\pars{Responsorium 3.} \scriptura{\Rbardot{} Prv. 8, 23-25 \Vbardot{} ibid. 8, 26.30; \textbf{H398}}

\vspace{-5mm}

\responsorium{I}{temporalia/resp-inprincipiodeusantequam-CROCHU-cumdox.gtex}{}}
\newcommand{\hymnuslaudes}{\pars{Hymnus}

\grechangedim{interwordspacetext}{0.16 cm plus 0.15 cm minus 0.05 cm}{scalable}%
\cuminitiali{I}{temporalia/hym-DeusQui-hk.gtex}
\grechangedim{interwordspacetext}{0.22 cm plus 0.15 cm minus 0.05 cm}{scalable}%
}
\include{hebdomadaxviii}
\include{feriavi}
